The semantic grouping here is \term{expressive signs}: conventional uses
whose deployment conveys an evaluative or affective stance, rather than
merely describing that stance. This intended contribution is narrower than
all non-descriptive or use-conditional meaning. A conventional instruction
for managing interpretation or interaction needn't express such a stance.
Potts's diagnostics include perspective dependence, immediacy, and
intensifying repetition \citep[166--167]{potts2007}. For \mention{damn},
these ask whose attitude is conveyed, whether the utterance displays
that attitude on the occasion of speaking, and whether repeating the
form strengthens it. Whether those properties travel together across
the proposed extension is an empirical question.

\subsection{Contribution, encoding, and competing analyses}\label{sec:encoding}

Consider the constructed utterance \mention{The damn car won't start},
said after several failed attempts at ignition. The speaker reports a
starting problem and conveys a negative attitude. Calling \mention{damn}
expressive identifies that second contribution. Asking how the wording
helps a hearer recover the attitude raises a different question about
its encoding.

\textcite{wharton2016expressives} develops a procedural account of expressive
meaning, connecting guidance for interpretation with the showing of emotion.
Applied to the car example, such an account would treat \mention{damn}
as a cue for recognizing the attitude being shown. The use could then
be expressive in contribution and procedural in encoding. Those descriptions
would answer different questions about the same utterance.

Now imagine a narrator representing the car's backfire with \mention{Bang!}.
The sign depicts a noise, and its performance can resemble what it
depicts. That iconic presentation differs from \mention{damn}'s conveyed
stance. Identifying the depiction still leaves open how the sign guides
interpretation: resemblance alone doesn't decide whether it supplies
conceptual content, an interpretive procedure, or some combination.

Two accounts would compete if they made exclusive claims
about the same \mention{damn} use. Suppose one assigned it only
conceptual content, such as a negative evaluation, while the other
assigned it only a procedure for recognizing an attitude. These are
hypothetical rivals, not summaries of Potts and Wharton. The first
needn't make annoyance part of what the sentence asserts. Their
incompatible claims about encoding would require adjudication.

Finer groupings also become visible in ordinary situations. Someone
touches a hot pan and says \mention{Ouch!}; someone finds missing keys
under a newspaper and says \mention{Aha!}. The situations are constructed;
\textcite[113]{ameka1992} gives the forms as examples of emotive and
cognitive interjections. The first gives voice to a sensation, the
second displays a change in knowledge or recognition.

Ameka's mental-state grouping is broader than the stance grouping used
here. Finding the keys can be displayed with \mention{Aha!} without
conventionally evaluating them. That reading motivates a cognitive subtype;
it doesn't prevent another occurrence from also conveying affect.
Different contributions, crosscutting accounts of encoding, and rival
encoding hypotheses thus have different consequences even within this
small set of expressions.

Consider two constructed exchanges. A teller says \mention{I tried the
front door}, a listener responds \mention{Mm-hm}, and the teller continues
\mention{but it was locked}. The listener supports continuation of the
story. In a second exchange, someone says \mention{We meet at five};
a recipient who missed the time responds \mention{Huh?}, and the first
speaker repeats \mention{At five}. Here the response prompts repair.

The difference concerns what happens next: the story proceeds after
the continuer, while the meeting time is repeated after the repair
initiator. Closing \mention{bye} has another job, participating in the
end of an encounter. Accounts of interjections in conversation make
these sequential differences central \citep[257--259, 263--265]{dingemanse2024}.
Their union doesn't supply one expectation about the next turn.

The \term{bearer} is the object grouped. Two recordings of \mention{The
damn car won't start} contain two realizations of the selected prenominal
\mention{damn} use. Counting the recordings gives two occurrences, not
two sign-types. A phonologist can classify the individual productions;
a lexical inquiry can ask about the conventional use they realize.
A historian may instead classify a change connecting an earlier sign
with a later one.

The \term{grouping basis} supplies the assignment evidence. Prenominal
\mention{damn} enters the expressive grouping because of its conveyed
stance, while its position inside \mention{the damn car} supplies
grammatical evidence. Grouping it with standalone \mention{Damn!} for
the first reason doesn't erase their difference in the second respect.
Other proposals can use ancestry or participation in a social repertoire;
each needs to identify the relation its membership records.

Three constructed comparisons distinguish the reach of an expectation.
After identifying the expressive contribution in \mention{The damn car
won't start}, an investigator could examine another speaker's utterance
of that sentence, the question \mention{Has the damn car broken down?},
or \mention{The bloody car won't start}. These comparisons change the
realization, the environment, and the expression respectively.

The question supplies a concrete target: does the negative stance
remain conveyed while the car's breakdown is put in question? Evidence
from declaratives needn't settle that. The investigator can test the
familiar property in a new environment without inventing a differently
named property. Generalizing to independently identified \mention{bloody}
makes a further commitment across sign-types. Neither move is licensed
simply by defining the class to contain the desired cases.

Imagine an investigator who's identified expressive \mention{damn}
only in standalone reactions such as \mention{Damn! We've missed the
train}. Their provisional identification uses the reaction meaning;
it doesn't require standalone position. They predict that these uses
resist grammatical integration. Encountering \mention{The damn car won't
start} then presents an adverse case: the expression conveys stance
while functioning as an integrated modifier inside the noun phrase.

This is a hypothetical research sequence using familiar constructions.
The investigator should check the structural analysis, but the confirmed
prenominal occurrence counts against the restriction under the original
identification rule. Its inconvenient distribution isn't by itself a
reason to remove it from that assessment. Later recognizing a distinct
prenominal use can improve the analysis without making the original
prediction successful.

Suppose an inventory records that a use occurs standalone and conveys
stance. Adding the label \enquote{standalone stance use} abbreviates
those two observations; it adds no information to a fully observed
record. The combined grouping could still be useful when investigating
something not yet recorded, such as interpretation under repetition.
Pooling observations from its members might help estimate that pattern
from sparse records. Whether it does is a comparative question.

The grouping selects an intersection: stance \mention{Damn!} qualifies,
while prenominal \mention{damn} in \mention{the damn car} fails the
standalone condition. A model can represent that intersection as one
label or keep its two component assignments. Their relative value
depends on the information available and the target being predicted.
A benefit for estimating repetition behavior wouldn't establish a
unique cause of the shared properties or priority for every inquiry.

Explanation asks another question. Once the present form of \mention{gee}
has been recorded, knowing its religious source adds nothing to that
already observed spelling. Its history can still explain how the form
arose. More generally, conditioning a prediction on consequences of a
historical process can remove the association through which that process
is recognized. Section~\ref{sec:history} distinguishes this predictive
comparison from assessment of the proposed history itself.

An account of the stratum would gain projective reach if membership
organized acquisition, adaptation, or alternation beyond the diagnostic
used to identify it. Consider the contrast between producing a click
and pronouncing \mention{tut-tut} as spelled words. An investigation
could first group signs by independently recorded articulations, then
ask which permit a spelling pronunciation and in what settings.
That prospective comparison would test a relation beyond the articulation
used for grouping; the reported alternatives don't yet establish it.

Take a recording of \mention{Damn, we're going to be late!}. The
context can identify the reaction use while the acoustic record supplies
evidence about a boundary, such as a pause or pitch movement. A
second recording can realize the same use with different timing.
Classifying these productions asks how that occurrence was packaged;
a claim about characteristic packaging needs the distribution of
productions across relevant environments.

The comma in a transcript of \mention{Damn, we're going to be late!}
is a different kind of observation. It may record a transcriber's
analysis of detachment, but it isn't itself a pause or a grammatical
relation. Entering that comma once as a prosodic boundary and again
as syntactic independence would manufacture apparent agreement between
the measures. An acoustic analysis and a constructional analysis can
converge without relying on that single transcription decision.

Evidence for replication also needs to distinguish innovations from their
descendants. Suppose one community coins a clipped exclamation and
three others borrow it. Four community records then document one
origin, not four independent instances of clipping. They can support
an account of transmission or subsequent divergence. Testing whether
the same source repeatedly gives rise to the same outcome would
require independently established innovations. This hypothetical count
makes the sampling unit depend on the historical claim.

Consider a proposed comparison of \mention{Gee!} and \mention{Jesus!}
as reactions to surprising news among close friends and during a
formal interview. Ratings would ask how appropriate each choice seems
in each setting. A listener actually correcting a speaker's choice
would be a different observation. Repeated corrections followed by
changes in use would be relevant to a further claim about regulation.
The comparison specifies possible evidence, not findings about these forms.

For example, if listeners regularly associated \mention{Gee!} with an
old-fashioned speaker persona, their judgments would support a social
association. They wouldn't merely be rating whether it was acceptable
in the interview. Both investigations could concern the same surprise
reaction, while asking different questions about its social meaning.
Such associations could also unite forms used for several actions.
Any claimed convention would need a specified community and evidence
of shared recognition.

Compare a familiar greeting, \mention{How are you?}, a pause filler,
\mention{um}, and the idiom \mention{kick the bucket}. These illustrative
forms differ in grammar and conversational job. A processing inquiry
could ask whether speakers produce them as familiar routines.
Its grouping would cross the earlier grammatical and semantic divisions;
\term{formulaic expressions} might name it more clearly than another
kind of interjection.

For the greeting, one question is whether \mention{How are you?} is
retrieved fluently because repeated use has helped establish a routine.
A comparison at the same measured frequency asks how fluency differs
among expressions with comparable recorded use. That adjustment doesn't
by itself identify how much repetition contributed to the greeting's
fluency. A causal interpretation still needs assumptions about exposure,
retrieval, and how the frequency measure relates to speakers' experience.

Suppose retrieval across \mention{How are you?}, \mention{um}, and
\mention{kick the bucket} tracked familiarity more closely than any
interjection grouping. That result would favour the broader or continuous
representation for this processing inquiry. It wouldn't give the filler
\mention{um} the internal grammar of \mention{kick the bucket}. The
grammatical contrast would remain available for a different question,
provided its own evidence held.

\begin{table}[htbp]
\centering
\small
\caption{Further inquiries illustrated by constructed comparisons}
\label{tab:extensions}
\begin{tabular}{@{}>{\raggedright\arraybackslash}p{.19\textwidth}
                  >{\raggedright\arraybackslash}p{.31\textwidth}
                  >{\raggedright\arraybackslash}p{.42\textwidth}@{}}
\toprule
Inquiry & Possible bearer & Concrete question to assess \\
\midrule
Morphological organization & Sign-types in identified constructions &
Does exclamatory \mention{heavens!} participate in a productive number contrast? \\
\addlinespace
Semiotic relations & Conventional signs, or sign--referent relations &
How do stance \mention{Damn!} and depictive \mention{Bang!} relate form to interpretation? \\
\addlinespace
Development and learning & Learners' repertoires or acquisition episodes &
What changes between a child's first and later uses of \mention{ouch}? \\
\addlinespace
Multimodal organization & Conventional complexes or situated performances &
How does \mention{Bye!} coordinate with a wave, and what does each component contribute? \\
\addlinespace
Graphic communication & Written sign-types or particular inscriptions &
What does \mention{BANG} in a comic panel represent, compared with \mention{bang} in a transcript? \\
\addlinespace
Annotation & Tokens under a stated coding scheme &
Should a query return both \mention{Ha! I knew it} and \mention{12 ha of land}? \\
\bottomrule
\end{tabular}
\end{table}

In a constructed comic panel, \mention{BANG} beside a backfiring car
can depict its noise without representing anyone's speech. In a
transcript, \mention{bang} can instead record a narrator's spoken depiction.
The inscriptions have the same letters but stand in different relations
to speech and events. Work on ideophones distinguishes depiction from
the predominantly indexical relation associated with interjections
\citep{dingemanse2018,dingemanseakita2017}; identifying the relation in a
particular inscription needs evidence about its use.

Following a child learning \mention{ouch} would record a learner,
occasions of use, and changes between them. Those episodes aren't
additional entries in an adult inventory. Likewise, a recording of
\mention{Bye!} with a wave includes the timing of a visible movement,
which the lexical sign alone doesn't specify. Acquisition and performance
can be classified, but their connection to sign-types runs through
learning and realization rather than shared membership in one list.

The constructed strings \mention{12 ha of land} and \mention{Ha! I
knew it} make the difference visible. A search for the spelling
\mention{ha}, ignoring capitalization, can retrieve both. An investigation
of the reaction use needs to distinguish them. A tag intended to
make that distinction is still a measurement decision whose accuracy
needs checking.

In \mention{The damn car won't start}, expressive contribution concerns
the attitude conveyed; procedural encoding concerns how the wording
guides its recognition. Both descriptions could hold. The contrast with
depictive \mention{Bang!} concerns a different contribution, while the
exclusive encoding hypotheses in Section~\ref{sec:encoding} disagree about
the same use. Assigning the hypotheses to different fields wouldn't
make their claims compatible.

Failure has correspondingly local consequences. Suppose a repetition
test found that an independently identified expressive use didn't
strengthen stance under the specified conditions. The prediction of
intensification would fail there; the observed distinction between
conveying stance and depicting a sound wouldn't disappear with it.
Restricting the prediction to another subtype would create a new
claim. Its support couldn't be inherited from the failed generalization.

Selecting standalone stance uses includes expressive \mention{Damn!}
and excludes prenominal \mention{damn} in \mention{the damn car}. That
intersection is well-defined and can fix a useful comparison set.
Its further value depends on the question: does restricting attention
to those uses help with the proposed inference? Neither the overlap
itself nor the availability of its component assignments answers that.

For \mention{Bang!}, uncertainty between a nominal-fragment and an
interjection analysis doesn't itself establish graded membership. The
analysis may simply be unresolved. A separate inquiry could ask how
typical speakers judge it compared with \mention{Damn!}; prototype
accounts organize such typicality and similarity \citep{rosch1975,taylor2003}.
Either a graded or a categorical representation can serve a task,
but its form alone doesn't establish what it licenses an investigator
to infer.

The contrast between stance \mention{Damn!} and depictive \mention{Bang!}
shows why one semantic inquiry can need several groupings. Unusual
phonation in \mention{Bang!} wouldn't by itself decide its grammar.
The documented route to \mention{gee} explains an outcome without
validating an ancestry category. A query for \mention{ha} can aid
retrieval while also returning hectares.
In each case, making the bearer, grouping basis, and intended contribution
explicit changes what would count as support or counterevidence.
An undifferentiated category judgment conceals those differences; the
account developed here makes them available for adjudication.
